% RESULTADO_RECUPERADO: S15, 11c_gravedad_holonomia_rev6.tex:128--212;
% S01, 02_dinamica_tres_hojas.tex:155--219; pantalla y respuesta de30.
% Pruebas de curvatura, covariancia y Bianchi reunidas para autonomía.
% La realización suave conserva sus condiciones de holonomía explícitas.

\subsection{Realización geométrica del transporte y de la soldadura}
\label{vii:sec:realizacion-cartan}

La representación discreta del retorno y la soldadura de pantalla
determinan dos aspectos del transporte: la acumulación de memoria en
las rutas y la comparación de sus fibras. Su realización geométrica
conserva ambas operaciones. En esta sección se fijan el espacio de
llegada, la normalización dimensional de la conexión y las identidades
que utilizará su variación.

\subsubsection{Pantalla, cotétrada y conexión}

El cociente de pantalla \(Q_\square\), su forma lorentziana y la
soldadura \(\beta_m\) proceden de
\ref{vii:sec:pantalla-respuesta}. En cada celda, una realización de
pantalla \(J_{{\rm scr},m}\) transporta la forma y la soldadura
a la fibra geométrica correspondiente. Su imagen celular es
\[
 e_m=J_{{\rm scr},m}\beta_m.
\]
La escala \(\ell_m\) permanece dentro de la soldadura. Los
cambios de marco actúan sobre \(J_{{\rm scr},m}\), la conexión
y el transporte de las celdas conjuntamente.

Una realización suave de esos datos sobre una variedad orientada
\(\mathcal U\) de dimensión cuatro consiste en una cotétrada
no degenerada \(e=(e^I)\), una conexión de Lorentz
\(\omega=(\omega^I{}_J)\) y una realización de rutas
\(\mathfrak R_{\rm C}\) que satisfacen
\begin{equation}
 \operatorname{Hol}_{\mathcal A_{\ell_0}}
       (\mathfrak R_{\rm C}(\gamma))
 =\rho_{\rm C}(\operatorname{Hol}_{\rm TPK}(\gamma)).
 \label{vii:eq:compatibilidad-holonomia-geometrica}
\end{equation}
En esta igualdad matricial, la representación se normaliza mediante
\[
 N_{\ell_0}(R,a)=(\Lambda(R),a/\ell_0),\qquad
 \rho_{\rm C}=N_{\ell_0}\circ\rho_{\rm C}^{\rm disc},
\]
donde \(\Lambda:\operatorname{Spin}^+(3,1)\to SO^+(3,1)\)
es la representación vectorial. La aplicación \(N_{\ell_0}\)
preserva el producto semidirecto: dividir la traslación
\(a+\Lambda(R)b\) por \(\ell_0\) equivale a componer las
dos traslaciones normalizadas. Por ello la vuelta completa produce
\(72u\) en la conexión normalizada y \(72\ell_0u\) en la
coordenatización dimensional. El levantamiento de espín conserva, por separado,
\(R_4^{c_g(\gamma)}\); la representación vectorial tiene el
núcleo central \(\{1,-1\}\), de modo que se retiene ese
levantamiento al utilizar campos espinoriales.
Identidad, concatenación, inversión y subdivisión se mantienen en la
realización correspondiente. El emparejamiento de
pantalla y la condición de holonomía especifican los datos de la
realización utilizada en las ecuaciones de campo.

Con \(\eta=\operatorname{diag}(-1,1,1,1)\), se define
\(g=\eta_{IJ}e^I\otimes e^J\). Para una escala elemental
positiva \(\ell_0\), constante en la coordenatización, la conexión afín es
\begin{equation}
 \mathcal A_{\ell_0}
 =\begin{pmatrix}\omega&\ell_0^{-1}e\\0&0\end{pmatrix}.
 \label{vii:eq:conexion-afin-realizada}
\end{equation}
La cotétrada tiene dimensión de longitud; el factor
\(\ell_0^{-1}\) homogeneiza su componente traslacional.

\begin{proposition}[Curvatura y torsión de la conexión realizada]
\label{vii:prop:curvatura-afin}
La curvatura de \eqref{vii:eq:conexion-afin-realizada} es
\begin{equation}
 \mathcal F_{\ell_0}
 =\mathrm d\mathcal A_{\ell_0}
   +\mathcal A_{\ell_0}\wedge\mathcal A_{\ell_0}
 =\begin{pmatrix}F_\omega&\ell_0^{-1}T\\0&0\end{pmatrix},
 \quad F_\omega=\mathrm d\omega+\omega\wedge\omega,
 \quad T=\mathrm de+\omega\wedge e.
 \label{vii:eq:curvatura-afin-realizada}
\end{equation}
\end{proposition}
\begin{proof}
El producto de matrices de formas produce en el bloque diagonal
\(\omega\wedge\omega\) y en el traslacional
\(\ell_0^{-1}\omega\wedge e\). La derivada exterior aporta
\(\mathrm d\omega\) y \(\ell_0^{-1}\mathrm de\), pues
\(\ell_0\) es constante. La suma demuestra la igualdad.
\end{proof}

La curvatura describe el defecto de transporte de marcos y la torsión
el de la cotétrada. Ambas pertenecen a la misma conexión realizada;
sus componentes conservan dominios tensoriales distintos.

\begin{proposition}[Covariancia e identidades de Bianchi]
\label{vii:prop:bianchi-cartan}
Bajo un cambio de marco \(g_L:\mathcal U\to SO^+(3,1)\),
\[
 e'=g_Le,\qquad
 \omega'=g_L\omega g_L^{-1}-\mathrm dg_L\,g_L^{-1},
\]
se cumple \(T'=g_LT\) y \(F_{\omega'}=g_LF_\omega g_L^{-1}\).
Además,
\begin{equation}
 D_\omega T=F_\omega\wedge e,\qquad D_\omega F_\omega=0.
 \label{vii:eq:bianchi-realizacion}
\end{equation}
\end{proposition}
\begin{proof}
En \(\mathrm d(g_Le)+\omega'\wedge g_Le\), los términos
que contienen \(\mathrm dg_L\wedge e\) se cancelan.
La misma sustitución en la curvatura da su transformación por
conjugación. Para la primera identidad de Bianchi, se expande
\[
 D_\omega(\mathrm de+\omega\wedge e)
 =(\mathrm d\omega+\omega\wedge\omega)\wedge e.
\]
En \(\mathrm dF_\omega+\omega\wedge F_\omega-F_\omega\wedge\omega\),
los términos que contienen \(\mathrm d\omega\) y los productos
cúbicos se cancelan por pares. Esto prueba la segunda identidad.
\end{proof}

\subsubsection{Las tres realizaciones cuadráticas de curvatura}

La hoja \(\tau\in\{+1,0,-1\}\) del TRIT se representa en
el álgebra de Cartan mediante generadores \(J_{ab}=-J_{ba}\)
y \(P_a^{(\tau)}\), con relaciones
\begin{align*}
 [J_{ab},J_{cd}]&=
 \eta_{bc}J_{ad}-\eta_{ac}J_{bd}
 -\eta_{bd}J_{ac}+\eta_{ad}J_{bc},\\
 [J_{ab},P_c^{(\tau)}]&=
 \eta_{bc}P_a^{(\tau)}-\eta_{ac}P_b^{(\tau)},\\
 [P_a^{(\tau)},P_b^{(\tau)}]&=-\tau J_{ab}.
\end{align*}
Para \(\ell>0\) constante, se escribe
\[
 \mathcal A_\tau^{\rm Cart}
 =\tfrac12\omega^{ab}J_{ab}+\ell^{-1}e^aP_a^{(\tau)}.
\]

\begin{proposition}[Curvatura de la realización tritificada]
\label{vii:prop:curvatura-tritificada}
Con las relaciones anteriores,
\begin{equation}
 \mathcal F_\tau^{\rm Cart}
 =\tfrac12(F_\omega^{ab}-\tau\ell^{-2}e^a\wedge e^b)J_{ab}
     +\ell^{-1}T^aP_a^{(\tau)}.
 \label{vii:eq:curvatura-tritificada}
\end{equation}
\end{proposition}
\begin{proof}
Se expande \(\mathrm d\mathcal A+\tfrac12[\mathcal A,\mathcal A]\).
El corchete del sector de Lorentz produce \(F_\omega\), el
mixto produce \(D_\omega e\), y el traslacional aporta
\(-\tfrac12\tau\ell^{-2}e^a\wedge e^bJ_{ab}\).
\end{proof}

En la hoja central se recupera la conexión afín; las hojas extremas
conservan los dos signos del término de curvatura constante. La
representación mantiene así el régimen trítico de la ruta antes de
aplicar las ecuaciones variacionales. La igualdad de holonomía
\eqref{vii:eq:compatibilidad-holonomia-geometrica} corresponde a la
hoja central. En cada hoja extrema se utiliza la representación
\(\rho_{{\rm C},\tau}\) en su grupo de Cartan y su condición
de holonomía respectiva. La corriente y el tensor métrico utilizan
los campos y la representación de la misma realización.
